Food-insecurity decisions are made for local government areas (LGAs), whereas estimates are commonly reported for states. The number of areas above a policy threshold is a separate posterior quantity: its variance depends on pairwise exceedance probabilities, which marginal probabilities alone cannot provide. Local allocation therefore requires joint uncertainty about the count, not only a separate risk estimate for each area.

We analyse the 2024 Nigeria Demographic and Health Survey across 774 LGAs, where severe household food insecurity was 29.1\% nationally. A Bayesian hierarchical spatial model includes an unstructured area effect and an independently fitted cluster effect. At a 50\% prevalence threshold, the fitted expected count was 115 LGAs, with a 95\% interval of 101--130. An independence calculation would understate the count's standard deviation by 14\%; the marginals alone permit a standard deviation from 0.156 to 74.11, so the interval reflects the model's dependence structure rather than the marginals alone. The design leaves 120 LGAs without a sampled cluster, covering 11.04\% of the population. Consequently, the identified count spans 87--231, and 120 of the 144-area width rests on assumption.

With pass rules fixed before fitting, spatial cross-validation coverage was 0.890 at the exposure faced by unclustered LGAs and 0.775 when an entire region was withheld. These tests assess area-level prediction, not the count interval, which they do not validate. What transfers beyond this application are the sharp bounds and the pre-declared protocol, not a new estimator.
